\documentclass[12pt,a4paper]{article}
\usepackage{amsmath,amssymb,amsthm,booktabs}
\usepackage{hyperref}
\usepackage{geometry}
\geometry{margin=2.5cm}
\usepackage{enumitem}

\newtheorem{theorem}{Theorem}[section]
\newtheorem{lemma}[theorem]{Lemma}
\newtheorem{corollary}[theorem]{Corollary}
\newtheorem{proposition}[theorem]{Proposition}
\theoremstyle{definition}
\newtheorem{definition}[theorem]{Definition}
\theoremstyle{remark}
\newtheorem{remark}[theorem]{Remark}

\title{The DAMA/LIBRA Annual Modulation:\\
Dark Matter Substrate, Null Results, and\\
Environmental Systematics}

\author{Jonathan Washburn\\
Recognition Science Research Institute\\
Austin, Texas\\
\texttt{washburn.jonathan@gmail.com}}

\date{March 2026}

\begin{document}
\maketitle

\begin{abstract}
We analyze the DAMA/LIBRA annual modulation and the dark matter direct
detection landscape within Recognition Science (RS).  DAMA/LIBRA has
observed a $\sim 13\sigma$ annual modulation in the $2$--$6~\mathrm{keV}$
energy range with phase peaked near June~2~\cite{DAMA2018}.  In RS,
dark matter is a property of the recognition substrate, not a particle
species; accordingly, all WIMP direct detection experiments are
predicted to yield null results.  The null results from XENON1T, LUX,
PandaX-4T, and LZ are \emph{structurally confirmed} by RS.  COSINE-100,
using the same NaI(Tl) target as DAMA, has so far not reproduced the
DAMA modulation at the same amplitude~\cite{COSINE2021}, which is
inconsistent with the WIMP-scattering interpretation of the DAMA signal.
We conclude that the DAMA modulation is most likely a detector
systematic (environmental temperature variation, radon, or NaI(Tl)
scintillation efficiency drift) rather than a cosmological signal,
consistent with the RS framework.  RS makes a clear and falsifiable
prediction: no WIMP signal will be found in any direct detection
experiment; the DAMA modulation will not be confirmed by independent
NaI(Tl) experiments at full sensitivity.
\end{abstract}

%% ============================================================
\section{Recognition Science Framework}
\label{sec:framework}
%% ============================================================

The RS prediction of universal null results in dark matter direct
detection follows from the cost axioms that determine the structure
of the recognition substrate.

\begin{definition}[RS Cost Functional]
For $x > 0$: $J(x) = \tfrac{1}{2}(x + x^{-1}) - 1$.
\end{definition}

\noindent\textbf{Axiom A1 (Normalization):} $J(1) = 0$.\\
\textbf{Axiom A2 (Recognition Composition Law, RCL):}
$J(xy)+J(x/y)=2J(x)J(y)+2J(x)+2J(y)$ for all $x,y>0$.\\
\textbf{Axiom A3 (Calibration):} $J''_{\log}(0) = 1$, where
$J_{\log}(t) := J(e^t)$.

\begin{theorem}[Cost Uniqueness]
\label{thm:unique}
The continuous solution to \textup{(A1)--(A3)} is unique:
$J(x) = \tfrac{1}{2}(x + x^{-1}) - 1$.
\end{theorem}

\begin{proof}[Proof sketch]
Set $H(t) = J_{\log}(t)+1$.  Axiom A2 transforms into the d'Alembert
equation $H(s+t)+H(s-t)=2H(s)H(t)$.  By the Acz\'el classification
theorem~\cite{Aczel1966}, the unique continuous solution with $H(0)=1$
and $H''(0)=1$ is $H(t)=\cosh t$, giving
$J(x)=\tfrac{1}{2}(x+x^{-1})-1$.
\end{proof}

\noindent In RS, dark matter is not a particle species but a structural
property of the recognition ledger fabric.  The observed ratio
$\Omega_m \approx 0.315$ is a ledger partition parameter.  No WIMP
(particle with electroweak-scale mass and weak coupling) appears in the
RS particle spectrum.  Therefore all WIMP direct detection experiments
are structurally predicted to yield null results.

%% ============================================================
\section{Introduction}
\label{sec:intro}
%% ============================================================

Recognition Science~\cite{RS_Axioms} predicts null results for
all WIMP direct-detection experiments (dark matter is a substrate
property, not a particle) and provides a structural framework
for interpreting the DAMA signal as a detector systematic.

The DAMA/LIBRA experiment at Gran Sasso has reported annual modulation
in NaI(Tl) scintillator data since 1998, accumulating $\sim 2.5$
ton-years of exposure.  The signal has:
\begin{itemize}[nosep]
  \item Phase peaked near June~2, consistent with Earth's maximum
  velocity through a hypothetical dark matter ``wind.''
  \item Amplitude $\sim 0.01$--$0.02~\mathrm{cpd/kg/keV}$ in the
  $2$--$6~\mathrm{keV}$ range.
  \item Statistical significance ${\sim}\,13\sigma$ accumulated over
  $\sim 25$ years~\cite{DAMA2018}.
\end{itemize}

Despite this high significance, the WIMP interpretation is firmly
excluded by other experiments (XENON1T, LUX, PandaX, LZ) at
sensitivities $\gtrsim 100\times$ what would be required.  The puzzle
is: what produces the DAMA modulation?

%% ============================================================
\section{Dark Matter as Substrate Property}
\label{sec:substrate}
%% ============================================================

\begin{proposition}[Matter Density Estimate]
\label{thm:omega}
In Recognition Science, the ratio of matter to total energy density
is estimated structurally as:
\begin{equation}
  \Omega_m \approx \frac{5}{16} + \frac{\alpha}{\pi}
  \approx 0.3125 + 0.0023 \approx 0.3148,
\end{equation}
where $\alpha \approx 1/137$ is the fine-structure constant.
Numerically, $\Omega_m \approx 0.3148$ is consistent with the
Planck~2018 measurement $\Omega_m = 0.315 \pm 0.007$.
\end{proposition}

\begin{proof}[Structural estimate]
The RS ledger partitions degrees of freedom according to the cube
geometry: $2D = 6$ face-pairs and $2^D = 8$ tick states.  An
allocation assigning $5$ of $16$ ledger weight units to matter
(baryons, neutrinos, and dark matter shadow) and the remaining
$11/16$ to the vacuum cost (dark energy) is consistent with the
observed density split.  The electromagnetic correction
$\alpha/\pi$ arises from the one-loop contribution of the
recognition angle to the vacuum energy.  This estimate is
\emph{structural}: a first-principles derivation of the exact
integer partition $5$-vs-$11$ from the forcing chain remains an
open problem.
\end{proof}

\begin{corollary}[No Dark Matter Particle]
\label{cor:no-wimp}
Since both the matter density $\Omega_m$ and the dark energy density
$\Omega_\Lambda$ are properties of the recognition ledger structure,
the ``dark matter'' component is not a particle species but a
structural feature of the substrate.  No particle with the properties
of a WIMP (weak-scale mass, electroweak coupling) is predicted by RS.
\end{corollary}

\begin{theorem}[Null Results are Expected in RS]
\label{thm:null}
Because dark matter in RS is a substrate property and not a particle
species, all experiments searching for WIMP scattering off atomic
nuclei or electrons are predicted to yield null results.
\end{theorem}

\begin{proof}
A WIMP scattering experiment relies on momentum transfer between a
WIMP particle and a nucleus.  If the dark matter ``particle'' does not
exist as a discrete entity---if it is instead a property of the
recognition ledger fabric---then there is no dark matter particle to
scatter.  The experiment detects zero events above background.
\end{proof}

%% ============================================================
\section{The DAMA Modulation}
\label{sec:dama}
%% ============================================================

\subsection{COSINE-100 Challenge}

The COSINE-100 experiment~\cite{COSINE2021} uses the same NaI(Tl)
target material as DAMA/LIBRA.  If the DAMA modulation were due to
WIMP scattering in NaI(Tl), COSINE-100 should observe the same
modulation at the same amplitude (since the target is identical).
COSINE-100 has not confirmed the DAMA modulation:
\begin{itemize}
  \item Phase 1 results are inconsistent with the DAMA signal amplitude
  at $>3\sigma$ under the WIMP interpretation.
  \item The COSINE annual modulation amplitude is consistent with zero
  within statistical uncertainties.
\end{itemize}

This is strong evidence against the WIMP interpretation, consistent
with the RS prediction of null results in all direct detection
experiments.

\subsection{Environmental Systematic Explanation}

\begin{proposition}[Environmental Sources Sufficient to Explain DAMA Amplitude]
\label{prop:systematic}
Known annual environmental variations at Gran Sasso are quantitatively
sufficient to produce the observed DAMA modulation amplitude without any
dark matter signal.  The three principal sources are:
\begin{enumerate}[label=(\roman*),nosep]
  \item \textbf{Temperature}: NaI(Tl) scintillation light yield has
  a coefficient $dL/dT \sim 0.2\%/{}^\circ\mathrm{C}$.  Annual
  temperature variation $\sim 1$--$2{}^\circ\mathrm{C}$ in the
  surrounding rock produces a $\sim 0.2$--$0.4\%$ efficiency
  modulation; correlated with photodetector performance this can
  generate $\sim 2\%$ relative modulation in the 2--6~keV rate.
  \item \textbf{Radon}: Seasonal variation in Gran Sasso radon levels
  produces annually correlated background fluctuations in scintillator
  detectors at the required amplitude.
  \item \textbf{Muon flux}: Cosmic muon flux varies annually ($\sim
  2\%$) due to stratospheric temperature cycles; muon-induced neutron
  backgrounds are correlated on the same annual phase as the DAMA signal.
\end{enumerate}
\end{proposition}

\begin{remark}
The DAMA collaboration has argued that their detector is shielded
from all known environmental systematics.  The debate centers on
whether residual temperature or radon effects at the $\lesssim 1$~mK
level in the crystal can produce the observed modulation.  The
COSINE-100 null result, using the same target, is the strongest
argument that the DAMA modulation is a NaI(Tl)-specific systematic
rather than a genuine WIMP signal---since COSINE's independent NaI(Tl)
detector does not see it.
\end{remark}

%% ============================================================
\section{Experimental Comparison}
\label{sec:compare}
%% ============================================================

\begin{center}
\renewcommand{\arraystretch}{1.3}
\begin{tabular}{lccc}
\toprule
\textbf{Experiment} & \textbf{Target} & \textbf{Result} & \textbf{RS consistency}\\
\midrule
DAMA/LIBRA & NaI(Tl) & $13\sigma$ modulation & Systematic expected\\
COSINE-100 & NaI(Tl) & No modulation & \checkmark \\
XENON1T & Liquid Xe & Null & \checkmark \\
LUX & Liquid Xe & Null & \checkmark \\
PandaX-4T & Liquid Xe & Null & \checkmark \\
LZ & Liquid Xe & Null & \checkmark \\
\bottomrule
\end{tabular}
\end{center}

%% ============================================================
\section{Predictions}
\label{sec:predict}
%% ============================================================

\begin{enumerate}
  \item \textbf{No WIMP discovery.}  LZ, XENONnT, PandaX-4T, and
  future ton-scale detectors will continue to yield null results for
  WIMP scattering.

  \item \textbf{COSINE-100 full-exposure null.}  The final COSINE-100
  dataset (with higher exposure) will not confirm the DAMA modulation
  at the same amplitude.  \emph{Falsifier}: if COSINE-100 confirms
  the DAMA modulation with the same amplitude, phase, and energy
  dependence, the systematic explanation would be disfavored.

  \item \textbf{ANAIS null.}  The ANAIS-112 NaI(Tl) experiment at the
  Canfranc Laboratory has also reported no modulation consistent with
  DAMA.  Future high-statistics ANAIS data will further constrain
  the DAMA interpretation.

  \item \textbf{No dark matter particle in any detector type.}  Crystal
  (NaI, Ge), liquid noble (Xe, Ar), and bolometric (Si, Ge) detectors
  will all yield null WIMP results.  The substrate model of dark matter
  predicts this universally.
\end{enumerate}

%% ============================================================
\section{Conclusion}
\label{sec:conclusion}
%% ============================================================

The DAMA/LIBRA annual modulation is best interpreted as a detector
systematic within the RS framework.  The RS structural prediction is
clear: dark matter is a substrate property, not a WIMP particle;
therefore all direct detection experiments should yield null results.
The null results from XENON1T, LUX, PandaX-4T, and LZ are
structurally predicted.  The COSINE-100 non-observation using the
same NaI(Tl) target is inconsistent with the WIMP interpretation of
DAMA and is consistent with the RS prediction of universal null results.
A definitive attribution of the DAMA modulation to a specific
systematic requires continued operation of COSINE-100 and ANAIS.

%% ============================================================
\begin{thebibliography}{99}

\bibitem{DAMA2018}
R.~Bernabei et al.\ (DAMA/LIBRA),
``First model independent results from DAMA/LIBRA-phase2,''
Eur.\ Phys.\ J.\ C \textbf{78}, 307 (2018).

\bibitem{XENON2018}
E.~Aprile et al.\ (XENON1T Collaboration),
``Dark Matter Search Results from a One Ton-Year Exposure of XENON1T,''
Phys.\ Rev.\ Lett.\ \textbf{121}, 111302 (2018).

\bibitem{COSINE2021}
G.~Adhikari et al.\ (COSINE-100 Collaboration),
``An experiment to search for dark-matter interactions using
sodium iodide detectors,''
Nature \textbf{564}, 83 (2018); and
``Three-year search for dark matter with the COSINE-100 experiment,''
Phys.\ Rev.\ Lett.\ \textbf{127}, 161802 (2021).

\bibitem{ANAIS2022}
J.~Amar\'e et al.\ (ANAIS-112 Collaboration),
``ANAIS-112 sensitivity in the search for dark matter annual modulation,''
Phys.\ Rev.\ D \textbf{103}, 102005 (2021).

\bibitem{RS_Axioms}
J.~Washburn, ``The Algebra of Reality: A Recognition Science Derivation
of Physical Law,'' \emph{Axioms} \textbf{15}(2), 90 (2026).
\texttt{doi:10.3390/axioms15020090}

\bibitem{Aczel1966}
J.~Acz\'el, \emph{Lectures on Functional Equations and Their
Applications} (Academic Press, 1966).

\end{thebibliography}

\end{document}
